% =====================================================================
% Figure 4 -- Prediction-Isolated vs. Joint Formulation
%
% Side-by-side comparison of (a) the conventional prediction-isolated
% pipeline and (b) the joint formulation adopted in this paper.
% The only structural difference is the orange dashed feedback arc
% in panel (b) -- this carries the entire visual message.
%
% Standalone-compilable. Can also be \input from the main document.
% =====================================================================

\definecolor{Cinkblack}{HTML}{1A1A1A}
\definecolor{Crulegrey}{HTML}{B0B5BB}
\definecolor{Cbgblue}{HTML}{E3ECF7}
\definecolor{Cbgyellow}{HTML}{FBF4DC}
\definecolor{Cbggreen}{HTML}{E7F0E5}
\definecolor{Cbggrey}{HTML}{F2F4F7}
\definecolor{Caccorange}{HTML}{E8833A}
\definecolor{Caccnavy}{HTML}{1E3A8A}
\definecolor{Cevtolred}{HTML}{C0524A}
\definecolor{Claneborder}{HTML}{CBD2D9}

\begin{tikzpicture}[
		font=\sffamily\small,
		>=Stealth,
		every node/.append style={text=Cinkblack},
	]

	% ---------------- Styles ---------------- %
	\tikzset{
		mainbox/.style = {
				rectangle, rounded corners=3pt,
				draw=Cinkblack, line width=0.7pt,
				fill=white,
				text width=3.6cm,
				minimum height=1.4cm,
				inner sep=5pt,
				align=center,
				font=\sffamily\small,
			},
		predbox/.style = {
				mainbox, fill=Cbgblue,
			},
		planbox/.style = {
				mainbox, fill=Cbggreen,
			},
		metricbox/.style = {
				mainbox, fill=Cbggrey,
			},
		losslabel/.style = {
				rectangle, rounded corners=2pt,
				draw=Caccnavy, line width=0.5pt,
				fill=white,
				inner sep=3pt,
				font=\sffamily\footnotesize\itshape,
				text=Caccnavy,
			},
		losslabel/.append style = {align=center},
		jointloss/.style = {
				losslabel,
				draw=Caccorange,
				text=Caccorange,
				line width=0.8pt,
			},
		fwd/.style = {
				->, line width=0.7pt, draw=Cinkblack,
			},
		grad/.style = {
				->, line width=1.4pt, draw=Caccorange, dashed,
				dash pattern=on 5pt off 2.5pt,
			},
		paneltitle/.style = {
				font=\sffamily\large\bfseries,
				text=Cinkblack,
				align=center,
			},
		panelsub/.style = {
				font=\sffamily\footnotesize\itshape,
				text=black!55,
				align=center,
			},
		punchline/.style = {
				rectangle, rounded corners=2pt,
				inner sep=7pt,
				text width=6.5cm,
				align=center,
				font=\sffamily\footnotesize\itshape,
			},
	}

	% =====================================================================
	% PANEL (a): prediction-isolated formulation  (LEFT)
	% Three stacked boxes + L_ADE label on the side
	% =====================================================================

	% --- the three stacked boxes ---
	\node[predbox]   (a-pred)   at (0,0)         {Predictor $f_\theta$};
	\node[planbox]   (a-plan)   [below=8mm of a-pred] {Planner};
	\node[metricbox] (a-metric) [below=8mm of a-plan] {Operational metrics};

	% --- the loss label on the upper-left, pointing into the predictor ---
	\node[losslabel] (a-loss) at ($(a-pred.west)+(-2.5cm,0)$)
	{Loss:\\$\mathcal{L}_{\mathrm{ADE}}$};

	\draw[fwd] (a-loss.east) -- (a-pred.west);

	% --- forward arrows between the three boxes ---
	\draw[fwd] (a-pred.south)   -- (a-plan.north)
	node[midway, fill=white, inner sep=1pt, font=\sffamily\scriptsize\itshape,
		text=Caccnavy] {predictions};
	\draw[fwd] (a-plan.south)   -- (a-metric.north)
	node[midway, fill=white, inner sep=1pt, font=\sffamily\scriptsize\itshape,
		text=Caccnavy] {control actions};

	% --- panel title and subtitle ---
	\node[paneltitle, anchor=south]
	at ([yshift=18pt]a-pred.north)
	{(a) Prediction-Isolated};
	\node[panelsub, anchor=south]
	at ([yshift=2pt]a-pred.north)
	{conventional pipeline};

	% --- punch line below panel (a) ---
	\node[punchline, fill=Cbggrey, draw=Claneborder, line width=0.4pt,
	anchor=north]
	at ([yshift=-12mm]a-metric.south)
	{Predictor trained against displacement;\\
	planner adapts to whatever it receives.\\[2pt]
	{\footnotesize\bfseries No gradient flows back.}};

	% =====================================================================
	% PANEL (b): joint formulation  (RIGHT)
	% Same three boxes, but with an orange feedback arc
	% =====================================================================

	% --- horizontal offset between panels ---
	% Need enough room for: (a)'s right edge + central "vs." + (b)'s loss
	% label on the left. Increased from 6.5cm to give the divider real
	% breathing room.
	\def\panelgap{9cm}

	% --- the three stacked boxes ---
	\node[predbox]   (b-pred)   at ($(a-pred)+(\panelgap,0)$)   {Predictor $f_\theta$};
	\node[planbox]   (b-plan)   [below=8mm of b-pred] {Planner\\
		\scriptsize\itshape (differentiable)};
	\node[metricbox] (b-metric) [below=8mm of b-plan] {Operational metrics};

	\node[jointloss] (b-loss) at ($(b-pred.west)+(-2.5cm,0)$)
	{Loss:\\$\mathcal{L}_{\mathrm{TASL}}$};

		\draw[fwd] (b-loss.east) -- (b-pred.west);

		% --- forward arrows between the three boxes ---
		\draw[fwd] (b-pred.south)   -- (b-plan.north)
		node[midway, fill=white, inner sep=1pt, font=\sffamily\scriptsize\itshape,
			text=Caccnavy] {predictions};
		\draw[fwd] (b-plan.south)   -- (b-metric.north)
		node[midway, fill=white, inner sep=1pt, font=\sffamily\scriptsize\itshape,
			text=Caccnavy] {control actions};

		% --- THE KEY DIFFERENCE: orange dashed feedback arc, from metrics box
		%     curving up the right side back into the predictor box ---
		\coordinate (arc-start) at (b-metric.east);
		\coordinate (arc-cp1)   at ($(b-metric.east)+(2.2cm,0)$);
		\coordinate (arc-cp2)   at ($(b-pred.east)+(2.2cm,0)$);
		\coordinate (arc-end)   at (b-pred.east);

		\draw[grad]
		(arc-start) .. controls (arc-cp1) and (arc-cp2) .. (arc-end);

		\node[font=\sffamily\footnotesize\bfseries\itshape, text=Caccorange,
		anchor=west]
		at ($(arc-cp1)!0.5!(arc-cp2)+(0.1cm,0)$)
		{gradient of $\mathcal{L}_{\mathrm{TASL}}$};
	\node[font=\sffamily\scriptsize\itshape, text=Caccorange,
		anchor=west]
	at ($(arc-cp1)!0.5!(arc-cp2)+(0.1cm,-0.4cm)$)
	{through KKT \& $s_\theta$};

	% --- panel title and subtitle ---
	\node[paneltitle, anchor=south]
	at ([yshift=18pt]b-pred.north)
	{(b) Joint Formulation};
	\node[panelsub, anchor=south]
	at ([yshift=2pt]b-pred.north)
	{\textbf{this paper}};

	% --- punch line below panel (b) ---
	\node[punchline, fill=Cbgyellow, draw=Caccorange, line width=0.6pt,
	anchor=north]
	at ([yshift=-12mm]b-metric.south)
	{Operational signals \textbf{supervise the predictor}.\\[2pt]
	{\footnotesize Predictor learns where uncertainty matters most\\
	for downstream safety.}};

	% =====================================================================
	% Central divider with "vs." marker
	% =====================================================================
	\coordinate (div-top) at ($(a-pred.north)!0.5!(b-pred.north)+(0,1.0cm)$);
	\coordinate (div-bot) at ($(a-metric.south)!0.5!(b-metric.south)+(0,-1.5cm)$);

	\draw[Claneborder, line width=0.4pt, dotted]
	(div-top) -- (div-bot);

	\node[circle, draw=Cinkblack, fill=white, line width=0.5pt,
		inner sep=2pt, font=\sffamily\footnotesize\bfseries]
	at ($(a-pred)!0.5!(b-pred)+(0,-0.7cm)$) {vs.};

\end{tikzpicture}
